\section{Lista delle modifiche pre-implementazione}

A seguito delle valutazioni condotte e delle iterazioni effettuate sui mockup cliccabili, abbiamo definito una lista ristretta di modifiche e nuove funzionalità da implementare sul prototipo prima di procedere alla stesura del codice della webapp finale:

\begin{table}[H]
    \centering
    \small
    \renewcommand{\arraystretch}{1.3}
    \begin{tabular}{
        p{1.5cm}
        |>{\raggedright\arraybackslash}p{6.8cm}
        |>{\centering\arraybackslash}p{2.1cm}
        |>{\centering\arraybackslash}p{2.1cm}
    }
    \toprule
    \rowcolor{culturiaRed}
    \textcolor{culturiaWhite}{\textbf{ID}} &
    \textcolor{culturiaWhite}{\textbf{Descrizione Intervento}} &
    \textcolor{culturiaWhite}{\textbf{Ambito}} &
    \textcolor{culturiaWhite}{\textbf{Priorità}} \\
    \midrule
    \rowcolor{culturiaBeige}
    \textbf{INT-01} & \textbf{Realtà Aumentata per singolo monumento}:\newline Abilitare la scansione e l'esplorazione AR per un singolo monumento/opera d'arte isolato, svincolando l'utente dall'obbligo di avviare un intero tour tematico a più tappe. & App Mobile & \textbf{ALTA} \\
    
    \textbf{INT-02} & \textbf{Integrazione Cicerone IA in AR}:\newline Integrare il chatbot dell'assistente virtuale IA direttamente all'interno della schermata della fotocamera AR, consentendo di dialogare con l'IA senza uscire dalla visuale immersiva. & App Mobile & \textbf{MEDIA} \\
    
    \rowcolor{culturiaBeige}
    \textbf{INT-03} & \textbf{Pannello Lingua dei Segni (LIS / ASL)}:\newline Progettare e prototipare un pannello dedicato all'accessibilità con un interprete virtuale in lingua dei segni (LIS/ASL) integrato direttamente nella vista AR delle opere. & App Mobile & \textbf{MEDIA} \\
    \bottomrule
    \end{tabular}
    \caption{Lista delle modifiche da effettuare prima dell'implementazione.}
    \label{tab:priorita_modifiche}
\end{table}
